\documentclass[12pt]{article}

\usepackage[a4paper, margin=2cm]{geometry}
\usepackage{fontspec}
\usepackage{setspace}
\usepackage{tabularx, booktabs}
\usepackage{graphicx}
\usepackage{amsmath}
\usepackage{minted}
\usepackage{float}
\usepackage{caption}

\usepackage{polyglossia}
\setdefaultlanguage{russian}
\setotherlanguages{english}

\usepackage{fontspec}
\setmainfont{NewComputerModern10}
\setsansfont{NewComputerModernSans10}
\setmonofont[
  ItalicFont={Monaspace Radon}
]{Monaspace Neon}

\newcolumntype{Y}{>{\centering\arraybackslash}X}

\begin{document}
\begin{titlepage}
\centering

{Министерство науки и высшего образования Российской Федерации}
\vskip 2mm
{\large\scshape Новосибирский государственный технический университет}
\vskip 2mm
{Кафедра прикладной математики}

\vfill
\vskip 3cm

{\scshape Практическая работа \textnumero 3}
\vskip 3mm
\begin{spacing}{1.7}
    {\huge\bf Метод конечных элементов для гармонических задач}
\end{spacing}
\vskip 5mm
{\scriptsize ФПМИ, ПМ-24}
\vskip 2mm
{\large\it Параскун И., Герасименко В.}

\vfill

{преподаватели}
\vskip 2mm
{\scshape Вагин Д. В., д.т.н., доцент}
\vskip 1mm
{\scshape Задорожный А. Г., к.т.н., доцент}
\vskip 1mm
{\scshape Леонович Д. А., ассистент}

\vfill

{Новосибирск, \the\year}

\end{titlepage}

\setcounter{page}{2}

\section{Постановка задачи}

\noindent Построить МКЭ для нелинейного уравнения гиперболического типа

\begin{align}
  \label{eq:1} -\nabla(\lambda\nabla u) 
    + \sigma\frac{\partial u}{\partial t}
    + \chi\frac{\partial^2u}{\partial t^2} = f,
\end{align}

\vskip 4mm
\noindent в котором правая часть $f$ представима в виде

\begin{align}
  f(x, y, z, t) = f^s(x, y, z)sin(\omega t) + f^c(x, y, z)cos(\omega t),
\end{align}

\vskip 4mm
\noindent с учётом следующих краевых условий:

\begin{align}
  \lambda\frac{\partial u^s}{\partial n} \bigg|_{S_2} = \theta^s,\ \  &
  \lambda\frac{\partial u^c}{\partial n} \bigg|_{S_2} = \theta^c \\
  \lambda\frac{\partial u^s}{\partial n} \bigg|_{S_3} + \beta(u^s|_{S_3} - u^s_{\beta}) = 0,\ \  &
  \lambda\frac{\partial u^c}{\partial n} \bigg|_{S_3} + \beta(u^c|_{S_3} - u^c_{\beta}) = 0.
\end{align}

\section{Теоретическая часть}

\noindent Решение уравнения \ref{eq:1} может быть представлено в следующем виде:

\begin{align}
  \label{eq:2} u(x, y, z, t) = u^s(x, y, z)sin(\omega t) + u^c(x, y, z)cos(\omega t).
\end{align}

\vskip 4mm
\noindent Подставляя \ref{eq:2} в \ref{eq:1} и группируя по $sin(\omega t)$ и $cos(\omega t)$,
  получаем систему уравнений

\begin{align}\label{eq:3}
  \begin{cases}
    -\nabla(\lambda\nabla u^s) - \omega\sigma u^c - \omega^2\chi u^s = f^s, \\
    -\nabla(\lambda\nabla u^c) + \omega\sigma u^s - \omega^2\chi u^c = f^c.
  \end{cases}
\end{align}

\vskip 4mm
\noindent Эквивалентная вариационная постановка в форме уравнения Галёркина для
  системы уравнений \ref{eq:3}:

\begin{align}
  \begin{cases}
    \begin{split}
      \int_{\Omega}(\lambda\nabla u^s \nabla v_0 & - \omega\sigma u^c v_0 - \omega^2\chi u^s v_0) d\Omega +
      \int_{S_3}\beta u^s v_0 dS \\
      & = \int_{\Omega}f^s v_0 d\Omega +
          \int_{S_2}\theta^s v_0 dS +
          \int_{S_3}\beta u^s_{\beta} v_0 dS,
    \end{split} \\
    \begin{split}
      \int_{\Omega}(\lambda\nabla u^c \nabla v_0 & - \omega\sigma u^s v_0 - \omega^2\chi u^c v_0) d\Omega +
      \int_{S_3}\beta u^c v_0 dS \\
      & = \int_{\Omega}f^c v_0 d\Omega +
          \int_{S_2}\theta^c v_0 dS +
          \int_{S_3}\beta u^c_{\beta} v_0 dS,
    \end{split}
  \end{cases}
\end{align}

\vskip 5mm
\noindent Раскладывая функции $u^s$, $u^c$ и $v_0$ по базису, переходим 
к конечноэлементной системе из $2n$ уравнений c $2n$ неизвестными:

\begin{align}
  \begin{cases}
    \begin{split}
      \sum_{j=1}^{n}\bigg(
        \int_{\Omega}(\lambda\nabla \psi_j \nabla \psi_i & - \omega^2\chi\psi_j\psi_i)d\Omega +
        \int_{S_3}\beta\psi_j\psi_i dS\bigg
      )q^s_j - \omega\sum_{j=1}^{n}\bigg(\int_{\Omega}\sigma\psi_j\psi_i d\Omega\bigg)q^c_j \\
      & = \int_{\Omega}f^s \psi_i d\Omega +
          \int_{S_2}\theta^s \psi_i dS +
          \int_{S_3}\beta u^s_{\beta} \psi_i dS,\ \ i=1..n
    \end{split} \\
    \begin{split}
      \sum_{j=1}^{n}\bigg(
        \int_{\Omega}(\lambda\nabla \psi_j \nabla \psi_i & - \omega^2\chi\psi_j\psi_i)d\Omega +
        \int_{S_3}\beta\psi_j\psi_i dS\bigg
      )q^c_j - \omega\sum_{j=1}^{n}\bigg(\int_{\Omega}\sigma\psi_j\psi_i d\Omega\bigg)q^s_j \\
      & = \int_{\Omega}f^c \psi_i d\Omega +
          \int_{S_2}\theta^c \psi_i dS +
          \int_{S_3}\beta u^c_{\beta} \psi_i dS\ \ i=1..n
    \end{split}
  \end{cases}
\end{align}

\vskip 5mm
\noindent Вид трилинейных базисных функций:

\begin{align*}
  & X_1(x) = \frac{x_{p+1}-x}{h_x}, & &X_2(x) = \frac{x - x_{p}}{h_x}, & &h_x = x_{p+1} - x_p, \\
  & Y_1(y) = \frac{y_{s+1}-y}{h_y}, & &Y_2(y) = \frac{y - y_{s}}{h_y}, & &h_y = y_{s+1} - y_s, \\
  & Z_1(z) = \frac{z_{r+1}-z}{h_z}, & &Z_2(z) = \frac{z - z_{r}}{h_z}, & &h_z = z_{r+1} - z_r.
\end{align*}

\begin{align}
  \psi_i = X_{\mu(i)}Y_{\nu(i)}Z_{\vartheta(i)}.
\end{align}

\section{Описание тестирования программы}

\begin{figure}[H]
  \centering
  \vspace{5mm}
  \includegraphics[scale=1]{pde/test/msh/m1/geo*.png}
  \vspace{5mm}
  \caption{Засчётная область}
\end{figure}

\subsection{Сходимость по пространству}

\begin{table}[H]
  \begin{tabularx}{\textwidth}{lYYYY}
  \toprule
    $j$ & \multicolumn{4}{c}{$E(|u - u^*|)$}                            \\
  \cmidrule{2-5}
        & 1             & 8             & 512           & 4096          \\
  \cmidrule{1-5}
    10	& 4.0576985E+01	& 1.0144246E+01	& 6.3401538E-01	& 1.5850385E-01 \\
    11	& 3.7887393E+01	& 9.4718482E+00	& 5.9199051E-01	& 1.4799763E-01 \\
    12	& 3.3687350E+01	& 8.4218375E+00	& 5.2636485E-01	& 1.3159121E-01 \\
    13	& 2.8144299E+01	& 7.0360748E+00	& 4.3975467E-01	& 1.0993867E-01 \\
    14	& 2.1479224E+01	& 5.3698060E+00	& 3.3561287E-01	& 8.3903218E-02 \\
    15	& 1.3957840E+01	& 3.4894599E+00	& 2.1809124E-01	& 5.4522811E-02 \\
  \bottomrule
  \end{tabularx}
  \caption{Сходимость по пространству для функции $u = x^2ysin(\omega t) + zcos(\omega t)$}
\end{table}

\begin{table}[H]
  \begin{tabularx}{\textwidth}{lYYYY}
  \toprule
    $j$ & \multicolumn{4}{c}{$E(|u - u^*|)$}                            \\
  \cmidrule{2-5}
        & 1             & 8             & 512           & 4096          \\
  \cmidrule{1-5}
    10	& 4.0863998E+03	& 1.6734571E+03	& 1.4834158E+02	& 3.7241988E+01 \\
    11	& 3.8086475E+03	& 1.5997733E+03	& 1.3937143E+02	& 3.5285550E+01 \\
    12	& 3.4192178E+03	& 1.4979800E+03	& 1.2570095E+02	& 3.2181796E+01 \\
    13	& 2.9254489E+03	& 1.3422111E+03	& 1.0884979E+02	& 2.7902891E+01 \\
    14	& 2.3624655E+03	& 1.1329324E+03	& 8.7662530E+01	& 2.2511778E+01 \\
    15	& 1.7052980E+03	& 8.7848724E+02	& 6.2997284E+01	& 1.6223191E+01 \\
    16	& 1.1349167E+03	& 5.8936116E+02	& 4.4129266E+01	& 1.1020596E+01 \\
    17	& 1.4934257E+03	& 6.8033462E+02	& 5.4344345E+01	& 1.4249712E+01 \\
    18	& 2.0312950E+03	& 9.3009093E+02	& 7.7813133E+01	& 2.0141262E+01 \\
    19	& 2.5029231E+03	& 1.1553165E+03	& 9.8254924E+01	& 2.5234910E+01 \\
    20	& 2.9625875E+03	& 1.3726951E+03	& 1.1583692E+02	& 2.9390432E+01 \\
  \bottomrule
  \end{tabularx}
  \caption{Сходимость по пространству для функции $u = e^xsin(\omega t)$}
\end{table}

\subsection{Сходимость по времени}

\begin{figure}[H]
  \centering
  \includegraphics[scale=0.7]{pde/test/hmc/c1/05.png}
  \caption{Значение функции $u = xysin(\omega t) + zcos(\omega t)$ в точке $(5, 5, 5)$, $\Delta t = 0.5$}
\end{figure}

\begin{figure}[H]
  \centering
  \includegraphics[scale=0.7]{pde/test/hmc/c1/025.png}
  \caption{Значение функции $u = xysin(\omega t) + zcos(\omega t)$ в точке $(5, 5, 5)$, $\Delta t = 0.25$}
\end{figure}

\begin{figure}[H]
  \centering
  \includegraphics[scale=0.7]{pde/test/hmc/c1/0125.png}
  \caption{Значение функции $u = xysin(\omega t) + zcos(\omega t)$ в точке $(5, 5, 5)$, $\Delta t = 0.125$}
\end{figure}

\section*{Вывод}

\begin{itemize}
  \item В отношении сходимости по пространству, метод решения параболических задач с выделением 
    гармонической составляющей по времени, как и следовало ожидать, обладает схожими свойствами c 
    соответствующим МКЭ для эллиптических задач. При этом, естественно, метод дает точное решение по времени.
  \item Среди методов решения параболических задач, основанных на аппроксимации решения по времени
    с помощью базисных полиномов Лагранжа, наибольшим соответствием точному решению, естественно, обладает
    четырехслойная неявная схема, которая может быть использована в качестве альтернативы исследуемому методу.
\end{itemize}

\section*{Приложение}
 
\fontsize{10}{10}\selectfont
\selectlanguage{english}
\inputminted{c}{pde/src/fem/hmc/fem.c}
\inputminted{c}{pde/src/fem/hmc/lin/slv.c}
\inputminted{c}{pde/src/fem/hmc/lin/asm.c}

\end{document}
